% BOOK-FILE-ID: f-78429a43ee32
\subsection{変わらない量を使う}\label{節：変わらない量を使う}
漸化式で項がどんどん変わっていっても,その裏で変わらない量が隠れていることがある.
その変わらない量を見つけ出し,\sectionref*{節：隣接3項型}に帰着させることを目標にしよう.

\subsubsection{例題：数項から候補を作る}
\[
  a_1=a_2=1,\qquad
  a_{n+2}=\frac{a_{n+1}^2+1}{a_n}\quad(n\ge1)
\]
を考える.まずは実験である.初めの数項は$1,1,2,5,13,\ldots$となる.
正の2項から次の項も正になるので,分母が0になることはない.
この数列をじっと眺めると,\,$2=3\cdot1-1$, $5=3\cdot2-1$, $13=3\cdot5-2$となっていて,\,$a_{n+2}=3a_{n+1}-a_n$と予想できる.
しかし,これはあくまで予想であり,証明が必要である.
$x=a_n,y=a_{n+1},z=a_{n+2}$とすれば$xz=y^2+1$で,\,$x,z$が対称に現れる.
$x+z=Ky$を目標にすると,係数の候補は
\[
  K=\frac{x+z}{y}=\frac{x^2+y^2+1}{xy}
\]
となる.この式の値$K$が$n$によらず一定なら,\,$a_{n+2}=Ka_{n+1}-a_n$という隣接3項型として解ける.そこで,実際に一定となるかを確かめよう.

\paragraphbf{変わらないことを示す}
作った量を
\BookMathLayout{\[
    I_n=\frac{a_n^2+a_{n+1}^2+1}{a_na_{n+1}}
  \]}{\[
    \begin{aligned}
      I_n=\frac{a_n^2+a_{n+1}^2+1}{a_na_{n+1}}
    \end{aligned}
  \]}
とする.\,$xz=y^2+1$より,
\[
  \begin{aligned}
    \frac{x^2+y^2+1}{xy}&=\frac{x+z}{y},\\
    \frac{y^2+z^2+1}{yz}&=\frac{x+z}{y}
  \end{aligned}
\]
なので$I_{n+1}=I_n$である.
$I_1=3$だから,すべての$n$について
\[
  a_n+a_{n+2}=3a_{n+1}
\]
を得る.初めの二項を保ったまま\bookterm{linear}{線形}の\bookterm{recurrence}{漸化式}へ帰着できた.

\paragraphbf{一般項へ戻す}
$\alpha=\frac{3+\sqrt5}{2}$, $\beta=\frac{3-\sqrt5}{2}$とすると,
\bookterm{characteristic}{特性方程式}$x^2-3x+1=0$から
\[
  a_n=\frac{\alpha^{n-1}+\beta^{n-1}}2
  -\frac{\alpha^{n-1}-\beta^{n-1}}{2\sqrt5}
\]
を得る.初めの2項はともに1であり,この漸化式では初めの2項から全項が決まるので,元の数列に一致する.
また,\,$a_{n+2}=3a_{n+1}-a_n$から整数性も帰納的に分かる.

\subsubsection{不変量から漸化式を求める}
添字を1つ進めても値が変わらない量を\textbf{\bookterm{invariant}{不変量}}という.
この例題では,保存される量から足し算と定数倍の\bookterm{linear}{線形}の関係を取り出した.\footnote{この\bookterm{linear}{線形}の式と元の式が全く同じ条件を表すと無条件に主張したわけではなく,
  元の数列について保存から導いた関係である.}

\begin{bookpoint}
  \textbf{POINT}\par
  実験から予想した関係を,添字を1つ進めても変わらない量として示し,初項でその値を決める.
\end{bookpoint}
ただし,不変量が見つかったからといって,一般項が求まるとは限らない.
初期値を変えたときの係数は\bookproblemref{practice-connections}{2}{接続演習問2},
整数方程式を保つ漸化式は\bookproblemref{practice-entrance}{9}{発展問9}へ.
付録の\sectionref*{節：演習を別の視点から}でも発展問9を扱う.
